\documentclass[11pt]{article}
\usepackage[a4paper,margin=25mm]{geometry}
\usepackage{amsmath,amssymb,amsthm}
\usepackage{booktabs}
\usepackage{array}
\usepackage[hidelinks]{hyperref}
\usepackage{microtype}

\newtheorem{lemma}{Lemma}
\newtheorem{corollary}[lemma]{Corollary}
\newtheorem{proposition}[lemma]{Proposition}
\theoremstyle{remark}
\newtheorem{remark}[lemma]{Remark}
\newtheorem{observation}[lemma]{Observation}

\newcommand{\rot}{\rho}
\newcommand{\Dfour}{D_4}
\newcommand{\Csym}{\mathrm{rot4}}
\newcommand{\Vsym}{\mathrm{dia2}}

\title{Balanced orbit defects in half-turn-symmetric\\ no-three-in-line configurations,\\
with new $2n$-point configurations for $n=36,\,37,\,39$}
\author{Alex Komang\thanks{Independent researcher, Bali, Indonesia.  Correspondence: see the repository page.
AI/tool-use disclosure: Section~\ref{sec:auth}.}}
\date{August 2026 --- draft v0.8}

\begin{document}
\maketitle

\begin{abstract}
A $2n$-point no-three-in-line configuration on the $n\times n$ grid with half-turn symmetry can be written as the union
of the orbits of a larger subgroup $H\le\Dfour$ (the quarter turns, or the two diagonal reflections) that it contains,
together with a set of half-turn pairs that are not $H$-orbits --- its \emph{orbit defect} relative to $H$.  We prove a
\emph{balance lemma}: the defect pairs, read as arcs on the row classes $\{i,n-1-i\}$, form a balanced directed
multigraph, hence a union of directed cycles.  With one further observation (two half-turn pairs on the same long
diagonal are collinear) this classifies the defects with at most three pairs: one pair must be a diagonal loop ---
Flammenkamp's pseudo-class rct4; two pairs are two loops on different diagonals or a directed $2$-cycle, which is an
orbit of the diagonal reflections; three pairs form a directed $3$-cycle or a loop plus a $2$-cycle.  We enumerate the
corresponding families exhaustively for small $n$ with an exact branch-and-bound program (validated against OEIS A000769
and the class counts of Flammenkamp's database): the $3$-cycle family for odd $n\le25$, $33$ and (in part) $37$, $39$,
$41$; the two-loop family for even $n\le28$ and $n=36$; the mixed families for $n\le27$ (all empty).  Among the
configurations found are: for $n=36$, the first located $2n$-point configurations whose symmetry group is exactly the
half-turn (exactly three inequivalent ones in the two-loop family); for $n=37$ (two) and $n=39$ (four), the first located configurations of that
symmetry group outside the rct4 pseudo-class (Flammenkamp's rct4 configurations of 1997/98 already have exactly
half-turn symmetry); and a further one for $n=33$.  ``First located'' refers to the public files of Flammenkamp's
database as of 2026-08-11 (confirmed by its maintainer, who has added the configurations to the database).  Configurations,
programs and sweep journals are public.
\end{abstract}

\section{Introduction}
The no-three-in-line problem (Dudeney, 1917) asks for the largest set of points of the $n\times n$ grid $\{0,\dots,n-1\}^2$
containing no three collinear points.  Since a line parallel to an axis contains at most two points, at most $2n$ points can be
placed, and $2n$-point configurations are the objects of interest.  Their existence for every $n$ is a well-known open problem;
a probabilistic heuristic (Guy--Kelly, corrected form) suggests that for large $n$ fewer than $2n$ points can be placed.
Computationally, $2n$-point configurations were known for all $n\le 46$ and for $n=48,50,52$ for decades
\cite{Flammenkamp92,Flammenkamp98}; in 2025--2026 constraint programming \cite{Prellberg26}, GPU branch-and-bound
\cite{Riley} and SAT solving extended this to all $n\le 70$ and to $n=72,74,76$ (see the database \cite{FlammenkampDB},
which also records, for each $n$, the number of configurations by their exact symmetry class under the dihedral group
$\Dfour$ of the square, and marks the entries for which no configuration is known).

This note concerns those entries rather than the record sizes.  Flammenkamp's table lists, for each $n$, the number of
inequivalent $2n$-point configurations whose stabilizer in $\Dfour$ is exactly a given subgroup: trivial (\emph{iden}), the
half-turn (\emph{rot2}), one diagonal reflection (\emph{dia1}), one axis reflection (\emph{ort1}), the quarter turns
(\emph{rot4}), both diagonal reflections (\emph{dia2}), both axis reflections (\emph{ort2}) and the full group (\emph{full}).
Inside the half-turn class he keeps, for odd $n$, the pseudo-class \emph{rct4} (``rot4 symmetry except on the long
diagonals'': quarter-turn orbits off the diagonals plus one half-turn pair on one diagonal), whose configurations have
stabilizer exactly the half-turn but are listed in a separate column \texttt{c}, the column \texttt{:} containing the
remaining rot2 configurations \cite{FlammenkampDB}.  The rct4 configurations were found by Flammenkamp around the turn of
1996/97 and published in \cite{Flammenkamp98}; the odd-$n$ construction of \cite{Prellberg26} is of the same type.  As of
2026-08-11 the column \texttt{:} was complete for $n\le31$; for $n=36$ both columns were empty, and for $n=37,39$ the column
\texttt{:} was empty while column \texttt{c} listed $21$ and $33$ configurations.

\paragraph{Results.}
(i)~A balance lemma (Section~\ref{sec:lemma}) describing which orbit defects a half-turn-symmetric configuration may have
relative to a larger subgroup $H\le\Dfour$: the defect pairs must form a balanced directed multigraph on the row classes.
Together with the collinearity of two pairs on one diagonal this classifies the defects with at most three pairs
(Proposition~\ref{prop:small}): a single pair is a diagonal loop (this is exactly the rct4 structure); two pairs are two
loops on different diagonals or a directed $2$-cycle, and a directed $2$-cycle is itself an orbit of the diagonal
reflections; three pairs form a directed $3$-cycle or a loop plus a $2$-cycle.
(ii)~Exhaustive enumerations of the resulting families for small $n$ (Section~\ref{sec:results}), with all sub-classes
listed and their journals kept: the $3$-cycle family, the two-loop family and the mixed families.
(iii)~New $2n$-point configurations with stabilizer exactly the half-turn: three inequivalent ones for $n=36$ (two-loop
family; both database entries were empty), two for $n=37$ and four for $n=39$ ($3$-cycle family; the first located
configurations of that stabilizer that are not of rct4 type), and one more for $n=33$; the configurations are listed in
the Appendix and were verified by independent programs.
(iv)~A description of the exact solver and its validation, measured growth of the cost of finding one $2n$-point
configuration in the classes rot4, rct4 and rot2, and a list of negative results with their exact scope.

\paragraph{Related work.}
The best known lower bound, $3n/2-o(n)$ points, is the modular-hyperbola construction of Hall, Jackson, Sudbery and Wild
\cite{HJSW75}; it has not been improved since 1975 (see \cite{KNS25,GGGKKP26} and Problem 72 of \cite{Green}).  For $k\ge3$
points per line the asymptotic answer $kn$ was obtained recently by probabilistic methods \cite{KNS25,GGGKKP26}, which
explicitly do not apply to $k=2$.  The probabilistic heuristic of Guy and Kelly \cite{GuyKelly}, in its corrected form
\cite{Voutier26,Prellberg26}, predicts that $2n$-point configurations become rare for large $n$; Prellberg notes that the same
heuristic wrongly predicts fewer than one symmetric configuration \cite{Prellberg26}, and Flammenkamp explains the abundance
of rotationally symmetric configurations by the fact that a pair of points blocks only $O(\log n)$ cells in the rotational
classes against $\Theta(n)$ in the axis-reflection classes \cite{FlammenkampDB}.  Machine-learning approaches (PatternBoost,
reinforcement learning) have been compared with integer programming on this problem and do not scale beyond $n\approx14$
\cite{RJPDP25}.  The abstract content of the balance lemma --- a zero-divergence condition and its cycle decomposition --- is
standard (see e.g.\ \cite{EKM} for balanced two-coloured degree sequences and alternating cycles); its application to
symmetric no-three-in-line configurations, and the resulting search families, we have not found in the literature.
Flammenkamp uses ``defect'' for the shortfall from $2n$ points and calls a no-three-in-line set with exactly $2n$ points a
\emph{solution}; to avoid a clash we say \emph{orbit defect}, and our ``$2n$-point configuration'' is his ``solution''.

\section{Symmetry classes and the balance lemma}\label{sec:lemma}
Let $m=n-1$ and let $\Dfour$ act on the grid; write $\rot(u,v)=(v,m-u)$ for the quarter turn, so $\rot^2(u,v)=(m-u,m-v)$,
$\sigma(u,v)=(v,u)$ for the reflection in the main diagonal and $\sigma'(u,v)=(m-v,m-u)=\rot^2\sigma(u,v)$ for the reflection
in the anti-diagonal.  A $2n$-point no-three-in-line configuration has exactly two points in every row and every column.

Fix a subgroup $H\le\Dfour$ containing $\rot^2$ and a configuration $C$ with $\rot^2C=C$.  Let $B$ be the union of all
$H$-orbits contained in $C$ (the \emph{$H$-base} of $C$; it is $H$-invariant, and it is determined by $C$ and $H$).  Since
$C$ is $\rot^2$-invariant and $\rot^2\in H$, the complement $D=C\setminus B$ is a union of half-turn pairs
$\{p,\rot^2p\}$, none of which is an $H$-orbit; $D$ is the \emph{orbit defect} of $C$ relative to $H$.  For
$a\in\{0,\dots,m\}$ let $r_X(a)$, $c_X(a)$ denote the numbers of points of a set $X$ in row $a$ and in column $a$.

\begin{lemma}[balance]\label{lem:balance}
Suppose that $H$ contains a motion $g$ with $g(\{\text{column }a\})=\{\text{row }a\}$ for every $a$ (the quarter turn
$\rot$ and the reflection $\sigma$ have this property; note that $\sigma'$ alone maps column $a$ onto row $m-a$, but
$\sigma'\in H$ implies $\sigma=\rot^2\sigma'\in H$).  Then every $H$-orbit $O$ satisfies $r_O(a)=c_O(a)$ for all $a$, and
consequently the orbit defect $D$ of a $2n$-point configuration satisfies $r_D(a)=c_D(a)$ for all $a$.
\end{lemma}
\begin{proof}
$g$ maps $O\cap\{\text{column }a\}$ bijectively onto $O\cap\{\text{row }a\}$ because $gO=O$; hence $r_O(a)=c_O(a)$.
Summing over the orbits in $B$ and using $r_C(a)=c_C(a)=2$ gives $r_D(a)=c_D(a)$.
\end{proof}

Interpret the pairs of $D$ as arcs of a directed multigraph on the \emph{row classes} $\bar x=\{x,m-x\}$ (for odd $n$ the
central class $\{m/2\}$ is a singleton): the pair $\{(x,y),(m-x,m-y)\}$ occupies rows $x,m-x$ and columns $y,m-y$, i.e.\ it
is an arc from $\bar x$ to $\bar y$ (a loop if $\bar x=\bar y$; for $\bar x\ne\bar y$ two different pairs can join the same
two classes).  A pair with $x=m/2$ has both points in the central row and is an arc out of the central class; a pair with
$y=m/2$ is an arc into it; the central \emph{cell} $(m/2,m/2)$ cannot belong to a $\rot^2$-invariant configuration with two
points per row, but the central class may well be used.  Then $r_D(\bar x)=2\cdot\mathrm{outdeg}(\bar x)$ and
$c_D(\bar x)=2\cdot\mathrm{indeg}(\bar x)$, and Lemma~\ref{lem:balance} says:

\begin{corollary}\label{cor:euler}
The orbit defect $D$ is a balanced directed multigraph on the row classes (in-degree $=$ out-degree at every class), hence
its arc set decomposes into directed cycles (loops and $2$-cycles allowed).  In particular a single defect pair must be a
loop, i.e.\ must lie on a long diagonal.  For $H=\langle\rot\rangle$ every $H$-orbit in $C$ has four points (the centre cell
cannot occur), so $2n=4a+2|D|$ and $|D|\equiv n\pmod 2$; for $H=\langle\sigma,\sigma'\rangle$ the diagonal $2$-orbits
$\{(i,i),(m-i,m-i)\}$, $\{(i,m-i),(m-i,i)\}$ are $H$-orbits, and if $b$ of them occur then $b+|D|\equiv n\pmod2$.
\end{corollary}

For $H=\langle\rot\rangle$ (rot4) an $H$-orbit off the diagonals is $\{(u,v),(v,m-u),(m-u,m-v),(m-v,u)\}$; for
$H=\langle\sigma,\sigma'\rangle$ (dia2) it is $\{(u,v),(v,u),(m-u,m-v),(m-v,m-u)\}$.  We call these the C4 base and the
V2 base.  The following proposition lists all orbit defects with at most three pairs; the families it names are then the
natural next steps beyond rct4.

\begin{proposition}[orbit defects with at most three pairs]\label{prop:small}
Let $C$ be a $2n$-point configuration with $\rot^2C=C$ and let $D$ be its orbit defect relative to $H=\langle\rot\rangle$.
\begin{enumerate}
\item[(a)] $|D|=1$: $D$ is one loop, i.e.\ one half-turn pair on a long diagonal ($n$ odd) --- the rct4 structure.
\item[(b)] $|D|=2$ ($n$ even): either two loops at different classes on different diagonals, $\{(x,x),(m-x,m-x)\}$ and
 $\{(y,m-y),(m-y,y)\}$ with $\bar y\ne\bar x$ (\textbf{two-loop family}); or a directed $2$-cycle $\bar x\to\bar y\to\bar x$,
 which is a dia2-orbit $\{(u,v),(v,u),(m-u,m-v),(m-v,m-u)\}$ with $\bar u=\bar x$, $\bar v=\bar y$ that is not a C4-orbit
 (\textbf{mixed family} C4 base + one V2 orbit).
\item[(c)] $|D|=3$ ($n$ odd): either a directed $3$-cycle on three distinct classes, $(x,y),(y',z),(z',x')$ with
 $\bar y'=\bar y$, $\bar z'=\bar z$, $\bar x'=\bar x$, one of the classes possibly the central one (\textbf{$3$-cycle family});
 or one loop plus a directed $2$-cycle, i.e.\ a diagonal pair plus a dia2-orbit that is not a C4-orbit (\textbf{mixed family
 with a loop}).
\end{enumerate}
The same list, with the roles of C4- and dia2-orbits exchanged, holds for $H=\langle\sigma,\sigma'\rangle$; there a directed
$2$-cycle is a C4-orbit that is not a dia2-orbit (\textbf{mixed family} V2 base + one C4 orbit), and a loop is a diagonal
$2$-orbit and hence never a defect.
\end{proposition}
\begin{proof}
By Corollary~\ref{cor:euler}, $D$ is a union of directed cycles on the classes.  Two loops on the same diagonal give four
collinear points, and two loops at the same class $\bar x$ on the two diagonals form the C4-orbit of $(x,x)$, which belongs to
the base by definition; this leaves (a) and the first case of (b), and excludes three loops.  A directed $2$-cycle
$\bar x\to\bar y\to\bar x$ consists of two pairs $\{(u,v),(m-u,m-v)\}$ and $\{(v',u'),(m-v',m-u')\}$ with $\bar u=\bar u'=\bar x$,
$\bar v=\bar v'=\bar y$; the second pair is the image of the first under $\sigma$ (if $(v',u')\in\{(v,u),(m-v,m-u)\}$) or
under $\rot$ (if $(v',u')\in\{(v,m-u),(m-v,u)\}$); in the second case the union is the C4-orbit of $(u,v)$ and belongs to the
base, in the first case it is the dia2-orbit of $(u,v)$.  For $|D|=3$ the balanced multigraphs are a $3$-cycle, a loop plus a
$2$-cycle, or three loops.  For the V2 base the argument is the same with $\rot$ and $\sigma$ exchanged.
\end{proof}

A configuration of one of these families is not $H$-symmetric (the defect breaks $H$); its stabilizer in $\Dfour$ is some
subgroup containing $\rot^2$ and not containing $H$, and it must be computed: for every configuration reported below the
stabilizer was determined by testing all eight motions, and it is exactly $\{1,\rot^2\}$ in every case, so the configurations
fall into the rot2 column of the database.  (For $n$ even, a two-loop configuration is ``rot4 symmetric except on the long
diagonals''; the database keeps rct4 as a pseudo-class for odd $n$ only, so these configurations belong to the column
\texttt{:}.)

\section{The solver}\label{sec:solver}
All searches were done with an exact branch-and-bound program (\texttt{there.c}, about 600 lines of C, no external
libraries; the family sweeps used the variant \texttt{there\_twisted\_experiment.c} with the two bases and the defect built
in).  Cells are $t=un+v$; a state keeps two bit-matrices, \emph{alive} by rows and by columns.  Every pair of chosen
points ``shades'' all grid cells of the line through them (walked with the primitive direction obtained from a gcd table);
shaded cells die, as do all cells of a row or column that already holds two points.  A symmetry class is imposed by choosing
whole orbits: a cell is a candidate only if its whole orbit is alive, and the ``orbit-alive'' bit-masks are derived from the
alive matrices by transposition and bit reversal (one word operation per group element).  Branching uses the most constrained
reading class (row or column with the fewest orbit-alive cells) and dies as soon as some class cannot be completed.  A failed
candidate is killed for its siblings, so that a search with a fixed symmetry class is complete: ``no configuration'' is a
computer-assisted proof of absence for that class (trusted program, no independent certificate).  With restarts and randomized
tie-breaking the same program is used to find one configuration quickly.  A \texttt{PAIRS} option fixes a set of half-turn
pairs (the orbit defect), kills their images under the base group, and searches the remaining orbits of the base; the same
option with a diagonal pair reproduces rct4 exactly.  Invalid \texttt{PAIRS}/\texttt{FIX} input is rejected.

\paragraph{Validation.}
For $n\le 10$ the program enumerates all $2n$-point configurations without symmetry: $1,2,11,32,50,132,380,368,1135$ labelled
configurations for $n=2,\dots,10$, which reduce to $1,1,4,5,11,22,57,51,156$ classes under $\Dfour$, in agreement with OEIS
A000769.  For symmetric classes we compared the number $N_H$ of labelled configurations invariant under a fixed subgroup $H$
with the class counts $c_K$ of \cite{FlammenkampDB}.  For a normal subgroup $H$ (rot2, rot4, dia2, ort2, full) all
$|\Dfour|/|K|$ labelled images of a configuration of class $K\supseteq H$ are $H$-invariant, so $N_H=\sum_{K\supseteq H}
c_K\,|\Dfour|/|K|$: rot2 for $n=8$: $36=4\cdot7+2\cdot4$; $n=10$: $67=4\cdot13+2\cdot6+2\cdot1+1$; dia2 for
$n=11,\dots,25$: $0,2,2,0,0,0,2,4=2\cdot(0,1,1,0,0,0,1,2)$.  For $H=\langle\sigma\rangle$ (dia1) only two of the four
images of a dia1 configuration have stabilizer $\langle\sigma\rangle$ (the other two have $\langle\sigma'\rangle$), so
$N_H=2c_{\mathrm{dia1}}+2c_{\mathrm{dia2}}+c_{\mathrm{full}}$: $n=13$: $20=2\cdot9+2\cdot1$; $n=15$: $28=2\cdot13+2\cdot1$.
For rct4 the program's mode places the diagonal pair on the main diagonal and therefore sees two of the four labelled images
of each rct4 configuration: $n=9,17,19,21$: $2,2,4,2=2\cdot(1,1,2,1)$ --- all in agreement.  Two independent
implementations (the two agents' programs) agree on every enumeration they both performed.  Every configuration reported here
was finally re-checked by an independent script that tests all $\binom{2n}{3}$ triples in exact integer arithmetic and
computes the stabilizer explicitly.

\section{Results}\label{sec:results}

\subsection{Exhaustive enumerations of the defect families}
The defect sets were enumerated up to $\Dfour$ (one representative per orbit of $\Dfour$ acting on defects, computed as
canonical forms of the six or four cells); for each representative the sub-class (defect $+$ base) was searched exhaustively.
Table~\ref{tab:threecycle} gives, for odd $n$, the number of inequivalent $3$-cycle defects --- $4\binom h3+\binom h2$ with
$h=(n-1)/2$, the second term being the cycles through the central class --- and the number of pairwise inequivalent
configurations found for the two bases; Table~\ref{tab:twoloop} gives the two-loop family for even $n$; the mixed families
are reported in the text.  Rows marked ``complete'' were swept entirely, every sub-class terminating with ``exhausted''
(or ``dead'', when the fixed defect already violates the law: e.g.\ $24$ of the $2240$ non-central sub-classes at $n=33$);
counts marked with $\ge$ come from sweeps that were still running when this draft was written, with the number of
sub-classes already swept stated.  \emph{Correction to an earlier draft:} the generator of the $3$-cycle lists used until
v0.6 omitted the cycles through the central class (an external audit pointed this out); the omitted sub-classes were swept
afterwards and produced two additional configurations at $n=13$ (V2 base) and $n=21$ (C4 base) --- both already in the
database, whose rot2 column is complete for $n\le31$ --- none at $n=15,\dots,25$, $33$ and $37$, and two new ones at $n=39$; the corrected counts are
given below, and all journals are in the repository (\texttt{logs/}).

\begin{table}[ht]\centering\small
\begin{tabular}{r r r r p{5.9cm}}
\toprule
$n$ & defects mod $\Dfour$ (central) & C4 base & V2 base & status of the sweep\\
\midrule
13 & 95 (15) & 5 & 3 & complete\\
15 & 161 (21) & 0 & 0 & complete\\
17 & 252 (28) & 0 & 1 & complete\\
19 & 372 (36) & 0 & 0 & complete\\
21 & 525 (45) & 1 & 0 & complete\\
23 & 715 (55) & 0 & 0 & complete\\
25 & 946 (66) & 0 & 0 & complete\\
33 & 2360 (120) & 1 & 0 & complete (both bases)\\
37 & 3417 (153) & 2 & 0 & complete (both bases)\\
39 & 4047 (171) & 4 & 0 & complete (both bases): C4 base 4 configurations (2 through the central class), V2 base none\\
41 & 4750 (190) & $\ge0$ & 0 & V2 base complete (none); C4 base: all 190 central sub-classes complete (none), 3396 of the 4750 sub-classes swept (none), rest in progress\\
45 & 6391 (231) & --- & $\ge0$ & V2 base: 400 sub-classes swept (none); C4 base not started\\
\bottomrule
\end{tabular}
\caption{The $3$-cycle family (odd $n$): C4- or V2-orbits plus a directed $3$-cycle of half-turn pairs; numbers of pairwise
inequivalent configurations.  A sub-class is one defect (one directed $3$-cycle up to $\Dfour$) together with one base;
``complete'' means every sub-class was searched exhaustively.}\label{tab:threecycle}
\end{table}

\begin{table}[ht]\centering\small
\begin{tabular}{r r p{9.5cm}}
\toprule
$n$ & sub-classes $\{x,y\}$ mod $\Dfour$ & inequivalent configurations\\
\midrule
14--22 & 21, 28, 36, 45, 55 & 0 (complete)\\
24 & 66 & 2 (complete)\\
26 & 78 & 1 (complete)\\
28 & 91 & 1 (complete)\\
36 & 153 & 3 (complete: three sub-classes contain a configuration, two labelled solutions each, i.e.\ one configuration
up to the diagonal reflections; the sub-class $\{1,4\}$ needed a rerun without time limit, $4.9\cdot10^7$ nodes)\\
\bottomrule
\end{tabular}
\caption{The two-loop family (even $n$): C4-orbits plus one half-turn loop on each long diagonal.  A sub-class is an
unordered pair of row classes $\{x,y\}$, $0\le x<y<n/2$ ($x$ = class of the main-diagonal loop, $y$ = class of the
anti-diagonal loop, or vice versa: the quarter turn interchanges the two roles), so there are $n(n-2)/8$ sub-classes.}
\label{tab:twoloop}
\end{table}

\paragraph{Mixed families.}  The mixed families of Proposition~\ref{prop:small} --- C4 base plus one dia2-orbit that is not a
C4-orbit (plus one diagonal loop for odd $n$), and V2 base plus one C4-orbit that is not a dia2-orbit --- were swept
exhaustively for all $9\le n\le27$ (both parities; e.g.\ $1210$ and $55$ sub-classes mod $\Dfour$ at $n=23$, $1956$ and
$78$ at $n=27$): no configuration exists in them.  We have no structural explanation.

\subsection{New configurations with exactly half-turn symmetry}
Table~\ref{tab:new} lists the new configurations; their coordinates and their encodings in the 90-character format of
\cite{FlammenkampDB} are given in the Appendix.  ``First located'' refers to the public files of the database on 2026-08-11 and to a search of the literature; the maintainer
has since confirmed that the \texttt{:} entries for $n=36,37,39$ were empty on that date and has added the configurations of
Table~\ref{tab:new} to the database.  For $n=36$ both database entries were empty; for $n=37$ and $39$ the
column \texttt{c} (rct4) was populated, so the configurations below are new representatives of exact half-turn symmetry
outside the rct4 structure, not the first configurations with that stabilizer; the $n=33$ configuration is new but the entry
was not empty.  All ten configurations were checked to be pairwise inequivalent (canonical forms under $\Dfour$).  The
second and third $n=39$ configurations have their defects through the central class --- they were found only after the
correction of the generator described above; the fourth was found when the exhaustive C4-base sweep of $n=39$ completed
(18 August 2026) and is not yet in the database.

\begin{table}[ht]\centering\small
\begin{tabular}{r l l}
\toprule
$n$ & family and defect & remark\\
\midrule
36 & two loops: $(6,6),(29,29)$ and $(5,30),(30,5)$ & first located exact half-turn (both cells empty)\\
36 & two loops: $(1,1),(34,34)$ and $(9,26),(26,9)$ & second, inequivalent\\
36 & two loops: $(1,1),(34,34)$ and $(7,28),(28,7)$ & third, inequivalent\\
37 & 3-cycle $(1,3),(3,31),(31,1)$, C4 base & first located outside rct4 (column \texttt{:})\\
37 & 3-cycle $(2,4),(4,20),(20,34)$, C4 base & second, inequivalent\\
39 & 3-cycle $(4,8),(8,20),(20,34)$, C4 base & first located outside rct4 (column \texttt{:})\\
39 & 3-cycle $(1,5),(5,19),(19,1)$ through the central class, C4 base & second, inequivalent\\
39 & 3-cycle $(4,7),(7,19),(19,4)$ through the central class, C4 base & third, inequivalent\\
39 & 3-cycle $(1,3),(35,14),(24,1)$, C4 base & fourth, inequivalent (found by the exhaustive C4-base sweep, 18 August 2026)\\
33 & 3-cycle $(2,3),(3,19),(19,2)$, C4 base & new (entry not empty)\\
\bottomrule
\end{tabular}
\caption{New $2n$-point configurations with stabilizer exactly $\{1,\rot^2\}$.}\label{tab:new}
\end{table}

\subsection{Cost of finding one configuration, and growth}
For orientation we report measured costs of finding one $2n$-point configuration on the $n\times n$ grid with the solver of
Section~\ref{sec:solver}, on two consumer machines (Apple M3 laptop, 8 cores; a 12-core x86 server; 8--12 seeds in parallel;
wall time to the first configuration):
rot4 (even $n$): $n=36$: 10--60\,s, $n=40$: 70\,s; rct4 (odd $n$): $n=41$: 806\,s, $n=43$: about 6.7 core-hours,
$n=45$: none within 3.3 core-hours; rot2: $n=24$: 10\,s, $n=27$: 137\,s.  These are measured ratios for named
instances, not a complexity law: the cost grows by roughly a factor $8$--$10$ per $+4$ in $n$ for rot4/rct4 and by a factor
of about $6$ per $+2$ for rot2, consistent with the exhaustive tree sizes (rot4: $\times20$ per $+4$; rot2 and dia1:
$\times9$--$10$ per $+2$) divided by the slowly growing numbers of configurations, and consistent with the growth reported for
CP-SAT in \cite{Prellberg26}.  Extrapolated --- with the obvious uncertainty of an extrapolation over more than $25$ steps ---
to the open sizes $71,73,75,77,\dots$ this suggests $10^5$--$10^7$ core-hours for our program; the record computations of 2026
(SAT and GPU search on clusters) are of that order or use better tools.  Sub-class sweeps of the defect families are cheap
because a fixed defect shades a large part of the grid: a sub-class of the $3$-cycle family at $n=37$--$41$ is exhausted in
seconds to a few minutes.

\subsection{Negative results, with their scope}
\begin{itemize}
\item \emph{Consequence of the lemma.}  A single defect pair off the diagonals is impossible (Corollary~\ref{cor:euler}); the
``rot4 except one pair anywhere'' family is empty, as our enumerations for $n\le 27$ had shown before the lemma was found.
\item \emph{Computer-assisted, exhaustive.}  The mixed families are empty for $9\le n\le27$ (above); the two-loop family is
empty for $14\le n\le22$; the $3$-cycle family is empty for $n=15,19,23,25$.
\item \emph{Bounded search, no conclusion about existence.}  Shifting a known rot4 configuration for $n=44$ into the
$46\times46$ grid (which preserves all its collinearities) and freeing $k$ of its $22$ orbits, no completion exists for
$k\le15$ (exhaustive for those $k$); local search on exactly $2n$ points (min-conflicts with tabu, simulated annealing) failed
from $n\approx12$--$28$ on with our implementations and budgets, although it produces large partial configurations easily
(e.g.\ 134 points on the $71\times71$ grid); belief-propagation decimation without backtracking failed from $n=8$ on.  These
are statements about the specified methods, not about the configurations.
\item \emph{Empirical observation.}  For the rot4/rct4 configurations of the database with $41\le n\le76$ (142
configurations) the largest subset lying on one line, axis-parallel parabola or shifted hyperbola over $\mathbb F_p$,
maximised over primes $11\le p\le61$, is on average $16\%$ of the points (never above $22\%$), \emph{less} than for random
sets with two points per row ($19\%$ on average): relative to these curve families and primes, the finite extremal
configurations do not reduce modulo a prime to a low-degree curve (cf.\ Green's remark on Problem 72 of \cite{Green}), while
the best asymptotic constructions do.  On the algebraic side, the modular-hyperbola construction of \cite{HJSW75} is optimal
within its own $2p\times2p$ box, unions of two hyperbolae gain only a bounded amount as far as computed, and for $H(1)\cup H(-1)$ a lawful
subset has at most $(11/3+o(1))(p-1)$ points (an explicit line cover); see the companion note \cite{Companion}.
\end{itemize}

\section{Reproducibility}
The programs, the sub-class lists, the sweep journals and the certified configurations are in the public repository
\url{https://github.com/iwasborninbali/saturation} (immutable tag \texttt{defects-v0.8}; checksums in
\texttt{docs/RELEASES.md}):\\
\texttt{there.c}; \texttt{kit71/there\_twisted\_experiment.c} and its even-$n$ variant \texttt{bench/there\_tw\_even.c}
(the two bases and the \texttt{PAIRS} defect built in; validated against full enumerations of the rot4 and dia2 classes);
the generators \texttt{bench/defects3.py} (all $3$-cycle defects mod $\Dfour$, central class included) and
\texttt{bench/mixed3.py}; the sweep drivers \texttt{bench/family\_sweep2.sh}, \texttt{central\_sweep.sh},
\texttt{mixed\_sweep.sh}, \texttt{twoloop\_exh.sh} (same directory); the finished journals under \texttt{logs/sweeps/} (one line
per sub-class with the solver's termination status; \texttt{BOOK} lines with every labelled solution); \texttt{web/stab.py}
for the exact class and the database encoding; \texttt{docs/SUBMISSION.md} for the configurations; \texttt{make check} for a
one-command test.  Typical commands: \texttt{ALL=1 PAIRS="1,3;3,31;31,1" ./there 37 2} enumerates the sub-class that
contains the first $n=37$ configuration (one configuration, $1.4\cdot10^7$ nodes); \texttt{PAIRS="6,6;5,30" ./there 36 2}
finds the first $n=36$ configuration in about 100\,s and \texttt{PAIRS="1,1;9,26" ./there 36 2} the second in about 30\,s
(wall time on the laptop); the sub-class $\{5,6\}$ is exhausted in about $7$ min and contains exactly two labelled
configurations, i.e.\ one up to the diagonal reflections, which preserve both loops.  Journals of the early sweeps run on the
first agent's machine (the first halves of the lists for $n=39,41,45$) are being added; the $n=33$ row was rerun in full with
the corrected generator on a cloud machine (\texttt{logs/sweeps/family3full\_*\_n33.txt}: $2360$ sub-classes per base, one
configuration), and the $n=37$ row was completed by sweeping the $153$ central sub-classes for both bases.  All exhaustive statements are
computer-assisted results of a trusted program (validated as in Section~\ref{sec:solver}); no independent certificates of
non-existence are produced.

\section{AI/tool-use disclosure and authorship statement}\label{sec:auth}
Following the arXiv policy on generative AI, the AI systems involved are not listed as authors; their role is disclosed here.
This work was carried out by autonomous AI agents (Claude Fable~5, Anthropic), running as independent solvers on separate
machines and communicating through the repository, under the direction of the author of record; a third agent of the same
kind worked mainly on the companion note \cite{Companion}.  The agents wrote the programs, designed and ran the searches,
found and proved the balance lemma (independently of each other, in two formulations), verified each other's results, and
drafted and revised this text, including the corrections after two external LLM-assisted audits commissioned by the author
of record; the author of record posed the problem, provided the computing environment, made all decisions about scope and
publication, and takes responsibility for the content; the extent of his independent verification of the proofs and of the
computations is recorded in \texttt{docs/HUMAN\_AUDIT.md} of the repository and will be stated precisely in any submitted
version.  We estimate that about nine tenths of the technical work was done by the AI systems; we state this plainly because
it is true.  All configurations reported here are verifiable by anyone with the scripts in the repository (every triple
checked in exact integer arithmetic).

\section*{Acknowledgements}
We thank Achim Flammenkamp for maintaining the no-three-in-line database and its class counts, without which the empty cells
would not have been visible, and Thomas Prellberg for the detailed account of the CP-SAT computations.

\begin{thebibliography}{9}
\bibitem{Flammenkamp92} A.~Flammenkamp, Progress in the no-three-in-line problem, \emph{J. Combin. Theory Ser. A} 60 (1992)
305--311.
\bibitem{Flammenkamp98} A.~Flammenkamp, Progress in the no-three-in-line problem, II, \emph{J. Combin. Theory Ser. A} 81 (1998)
108--113.
\bibitem{FlammenkampDB} A.~Flammenkamp, The no-three-in-line problem, database and tables,
\url{https://wwwhomes.uni-bielefeld.de/achim/no3in/readme.html} (accessed 2026-08-17; table state 2026-08-11; symmetry
remarks at \texttt{symmetry\_remarks.html}).
\bibitem{Prellberg26} T.~Prellberg, Constraint satisfaction programming for the no-three-in-line problem, arXiv:2602.07751 (2026).
\bibitem{Riley} M.~Riley, no-three-in-line (GPU branch-and-bound), \url{https://github.com/mvr/no-three-in-line}.
\bibitem{GuyKelly} R.~K.~Guy and P.~A.~Kelly, The no-three-in-line problem, \emph{Canad. Math. Bull.} 11 (1968) 527--531.
\bibitem{HJSW75} R.~R.~Hall, T.~H.~Jackson, A.~Sudbery and K.~Wild, Some advances in the no-three-in-line problem,
\emph{J. Combin. Theory Ser. A} 18 (1975) 336--341.
\bibitem{KNS25} B.~Kov\'acs, Z.~L.~Nagy and D.~R.~Szab\'o, Randomised algebraic constructions for the no-$(k+1)$-in-line problem,
arXiv:2508.07632 (2025); Advances in Combinatorics 2026:7.
\bibitem{GGGKKP26} A.~Ghosal, R.~Goenka, A.~Grebennikov, P.~Keevash, M.~Kwan and H.~T.~Pham, No-$(k+1)$-in-line problem for
$k\ge3$, arXiv:2607.05255 (2026).
\bibitem{Voutier26} P.~M.~Voutier, On the Guy--Kelly conjecture for the no-three-in-line problem, arXiv:2603.00215 (2026).
\bibitem{Green} B.~Green, 100 open problems, manuscript (version of December 2025), Problem 72,
\url{https://people.maths.ox.ac.uk/greenbj/papers/open-problems.pdf}.
\bibitem{RJPDP25} P.~Ramanathan, P.~D.~Joshi, T.~Prellberg, R.~A.~Dandekar and S.~Panat, Three methods, one problem: classical
and AI approaches to no-three-in-line, arXiv:2512.11469 (2025).
\bibitem{EKM} P.~L.~Erd\H{o}s, Z.~Kir\'aly and I.~Mikl\'os, On the swap-distances of different realizations of a graphical
degree sequence, arXiv:1205.2842 (2012).
\bibitem{Companion} A.~Komang, Extremal no-three-in-line subsets of a modular hyperbola in the Hall--Jackson--Sudbery--Wild
window, companion note (draft v0.8, tag \texttt{hjsw-note-v0.8}), \url{https://github.com/iwasborninbali/saturation}
(\texttt{paper/hjsw\_window.pdf}).
\end{thebibliography}

\appendix
\section{The configurations}
Rows and columns are numbered $0,\dots,n-1$ from the top-left corner; the encoding is the class character followed by two
column characters per row, with the 90-character alphabet of \cite{FlammenkampDB}:\\
\centerline{\texttt{0-9A-Za-z\#\$\%\&@?!()[]<>\{\}=*+|-/\textasciitilde\^{}\_:;,.}}

\paragraph{$n=36$, 72 points, stabilizer $\{1,\rho^2\}$.} Encoding:\\
{\small \texttt{:3NEHFSVZ3FPU6E2BNQ8I5GDS08JOT\allowbreak YVXDP9Y1QAM2416BGRZ7MJUHR9COXL\allowbreak T5AKW047KILCW}}\\
Rows $u:\,v_1,v_2$ (row $u$ has points at columns $v_1,v_2$):\\
{\small
  0:\,3,23\quad 1:\,14,17\quad 2:\,15,28\quad 3:\,31,35\quad 4:\,3,15\quad 5:\,25,30\quad 6:\,6,14\quad 7:\,2,11 \\
  8:\,23,26\quad 9:\,8,18\quad 10:\,5,16\quad 11:\,13,28\quad 12:\,0,8\quad 13:\,19,24\quad 14:\,29,34\quad 15:\,31,33 \\
  16:\,13,25\quad 17:\,9,34\quad 18:\,1,26\quad 19:\,10,22\quad 20:\,2,4\quad 21:\,1,6\quad 22:\,11,16\quad 23:\,27,35 \\
  24:\,7,22\quad 25:\,19,30\quad 26:\,17,27\quad 27:\,9,12\quad 28:\,24,33\quad 29:\,21,29\quad 30:\,5,10\quad 31:\,20,32 \\
  32:\,0,4\quad 33:\,7,20\quad 34:\,18,21\quad 35:\,12,32
}

\paragraph{$n=37$ (configuration 1), 74 points, stabilizer $\{1,\rho^2\}$.} Encoding:\\
{\small \texttt{:GJ39EP7VENDZJRSX7P6ZCF28IQ4VW\allowbreak YKQFa06COUa0LAG245WAISYLO1UBT3\allowbreak 89H1NDM5TBMRXHK}}\\
Rows $u:\,v_1,v_2$ (row $u$ has points at columns $v_1,v_2$):\\
{\small
  0:\,16,19\quad 1:\,3,9\quad 2:\,14,25\quad 3:\,7,31\quad 4:\,14,23\quad 5:\,13,35\quad 6:\,19,27\quad 7:\,28,33 \\
  8:\,7,25\quad 9:\,6,35\quad 10:\,12,15\quad 11:\,2,8\quad 12:\,18,26\quad 13:\,4,31\quad 14:\,32,34\quad 15:\,20,26 \\
  16:\,15,36\quad 17:\,0,6\quad 18:\,12,24\quad 19:\,30,36\quad 20:\,0,21\quad 21:\,10,16\quad 22:\,2,4\quad 23:\,5,32 \\
  24:\,10,18\quad 25:\,28,34\quad 26:\,21,24\quad 27:\,1,30\quad 28:\,11,29\quad 29:\,3,8\quad 30:\,9,17\quad 31:\,1,23 \\
  32:\,13,22\quad 33:\,5,29\quad 34:\,11,22\quad 35:\,27,33\quad 36:\,17,20
}

\paragraph{$n=37$ (configuration 2), 74 points, stabilizer $\{1,\rho^2\}$.} Encoding:\\
{\small \texttt{:FJ7E4PBGKO6U5VDZDH9RCF2X4QSTI\allowbreak ZQa2X0SEM8a3Y0A1I78AW3YLO9RJN1\allowbreak N5V6UCGKPBWMTHL}}\\
Rows $u:\,v_1,v_2$ (row $u$ has points at columns $v_1,v_2$):\\
{\small
  0:\,15,19\quad 1:\,7,14\quad 2:\,4,25\quad 3:\,11,16\quad 4:\,20,24\quad 5:\,6,30\quad 6:\,5,31\quad 7:\,13,35 \\
  8:\,13,17\quad 9:\,9,27\quad 10:\,12,15\quad 11:\,2,33\quad 12:\,4,26\quad 13:\,28,29\quad 14:\,18,35\quad 15:\,26,36 \\
  16:\,2,33\quad 17:\,0,28\quad 18:\,14,22\quad 19:\,8,36\quad 20:\,3,34\quad 21:\,0,10\quad 22:\,1,18\quad 23:\,7,8 \\
  24:\,10,32\quad 25:\,3,34\quad 26:\,21,24\quad 27:\,9,27\quad 28:\,19,23\quad 29:\,1,23\quad 30:\,5,31\quad 31:\,6,30 \\
  32:\,12,16\quad 33:\,20,25\quad 34:\,11,32\quad 35:\,22,29\quad 36:\,17,21
}

\paragraph{$n=39$, 78 points, stabilizer $\{1,\rho^2\}$.} Encoding:\\
{\small \texttt{:CGAWLU7M8CAN1FPZ2K9TXbOPYc7BB\allowbreak L5W3c2E4JIKJYOa0Z6XHRRV04DE159\allowbreak TIa3DNbFSQUGV8H6SMQ}}\\
Rows $u:\,v_1,v_2$ (row $u$ has points at columns $v_1,v_2$):\\
{\small
  0:\,12,16\quad 1:\,10,32\quad 2:\,21,30\quad 3:\,7,22\quad 4:\,8,12\quad 5:\,10,23\quad 6:\,1,15\quad 7:\,25,35 \\
  8:\,2,20\quad 9:\,9,29\quad 10:\,33,37\quad 11:\,24,25\quad 12:\,34,38\quad 13:\,7,11\quad 14:\,11,21\quad 15:\,5,32 \\
  16:\,3,38\quad 17:\,2,14\quad 18:\,4,19\quad 19:\,18,20\quad 20:\,19,34\quad 21:\,24,36\quad 22:\,0,35\quad 23:\,6,33 \\
  24:\,17,27\quad 25:\,27,31\quad 26:\,0,4\quad 27:\,13,14\quad 28:\,1,5\quad 29:\,9,29\quad 30:\,18,36\quad 31:\,3,13 \\
  32:\,23,37\quad 33:\,15,28\quad 34:\,26,30\quad 35:\,16,31\quad 36:\,8,17\quad 37:\,6,28\quad 38:\,22,26
}

\paragraph{$n=33$, 66 points, stabilizer $\{1,\rho^2\}$.} Encoding:\\
{\small \texttt{:MOAC3H9J6GJOBS9B05PT0VPQHV5UE\allowbreak I2C4SKUEI2R1F671W37RWLN4L8DGQD\allowbreak NFTKM8A}}\\
Rows $u:\,v_1,v_2$ (row $u$ has points at columns $v_1,v_2$):\\
{\small
  0:\,22,24\quad 1:\,10,12\quad 2:\,3,17\quad 3:\,9,19\quad 4:\,6,16\quad 5:\,19,24\quad 6:\,11,28\quad 7:\,9,11 \\
  8:\,0,5\quad 9:\,25,29\quad 10:\,0,31\quad 11:\,25,26\quad 12:\,17,31\quad 13:\,5,30\quad 14:\,14,18\quad 15:\,2,12 \\
  16:\,4,28\quad 17:\,20,30\quad 18:\,14,18\quad 19:\,2,27\quad 20:\,1,15\quad 21:\,6,7\quad 22:\,1,32\quad 23:\,3,7 \\
  24:\,27,32\quad 25:\,21,23\quad 26:\,4,21\quad 27:\,8,13\quad 28:\,16,26\quad 29:\,13,23\quad 30:\,15,29\quad 31:\,20,22 \\
  32:\,8,10
}


\end{document}
